\section{Tetrahedral trigonometry}
\label{sec:tetrahedron}

Let \(P_1,P_2,P_3,P_4\in\R^3\) be the vertices of a nondegenerate Euclidean
tetrahedron.  Let \(F_i\) be the face opposite \(P_i\), let \(\Delta_i>0\) be
its area, and let \(\theta_{ij}\in(0,\pi)\) be the interior dihedral angle
along the edge \(P_iP_j\).

\subsection{Oriented face-area vectors}

For each face choose the outward unit normal \(n_i\) and define its doubled
area vector
\begin{equation}
 N_i=2\Delta_i n_i.                                         \tag{11.1}
\end{equation}

\begin{lemma}[Closure of face-area vectors]
\label{lem:face-closure}
The four outward face-area vectors satisfy
\begin{equation}
 N_1+N_2+N_3+N_4=0.                                        \tag{11.2}
\end{equation}
\end{lemma}

\begin{proof}
Translate \(P_4\) to the origin and put
\(a=P_1-P_4\), \(b=P_2-P_4\), \(c=P_3-P_4\), with
\(\detm[a\ b\ c]>0\).  The inward doubled normals of the first three faces
are \(b\crossp c\), \(c\crossp a\), \(a\crossp b\).  The oriented normal of
the face \(P_1P_2P_3\) is
\begin{align*}
 (b-a)\crossp(c-a)
 &=b\crossp c+c\crossp a+a\crossp b.
\end{align*}
It points away from \(P_4\), so its inward normal is the negative of this sum.
Thus the four inward vectors sum to zero; negating all four proves (11.2) for
outward vectors.  This is a coordinate polynomial identity.
\end{proof}

Faces \(F_j\) and \(F_k\) meet along the edge joining the two remaining
vertices \(P_i,P_l\).  The angle between their outward normals is
\(\pi-\theta_{il}\).  Hence
\begin{equation}
 N_j\dotp N_k
 =4\Delta_j\Delta_k\cos(\pi-\theta_{il})
 =-4\Delta_j\Delta_k\cos\theta_{il}.                       \tag{11.3}
\end{equation}
This relation may itself be taken as the constructive definition of the
circle-valued dihedral angle; the positive sine comes from the normalized
cross product of the two face normals.

\subsection{The tetrahedral cosine law}

\begin{theorem}[Cosine-law analogue for a tetrahedron]
\label{thm:tetra-cosine}
For every permutation \(\{i,j,k,l\}=\{1,2,3,4\}\),
\begin{align}
 \Delta_i^2={}&\Delta_j^2+\Delta_k^2+\Delta_l^2\\
 &-2\bigl(
 \Delta_j\Delta_k\cos\theta_{il}
 +\Delta_j\Delta_l\cos\theta_{ik}
 +\Delta_k\Delta_l\cos\theta_{ij}
 \bigr).                                                    \tag{11.4}
\end{align}
\end{theorem}

\begin{proof}
By closure,
\(N_i=-(N_j+N_k+N_l)\).  Take squared norms:
\begin{align*}
 4\Delta_i^2
 &=4\Delta_j^2+4\Delta_k^2+4\Delta_l^2\\
 &\quad+2N_j\dotp N_k+2N_j\dotp N_l+2N_k\dotp N_l.
\end{align*}
Insert (11.3).  The pair \((j,k)\) leaves the edge \((i,l)\), the pair
\((j,l)\) leaves \((i,k)\), and the pair \((k,l)\) leaves \((i,j)\).
Division by \(4\) yields (11.4).
\end{proof}

This proof is the exact three-dimensional analogue of expanding
\(\norm{x+y}^2\) in the planar cosine law.  The ``side lengths'' are now face
areas and the relevant angles are supplements of the angles between outward
normals.

\begin{corollary}[Linear normal relation]
\label{cor:tetra-linear}
Dotting (11.2) with \(n_i\) gives the stronger first-degree relation
\begin{equation}
 \Delta_i=
 \Delta_j\cos\theta_{kl}+
 \Delta_k\cos\theta_{jl}+
 \Delta_l\cos\theta_{jk}.                                 \tag{11.5}
\end{equation}
\end{corollary}

For a regular tetrahedron, all four areas are equal and (11.5) gives
\(3\cos\theta=1\), the familiar \(\cos\theta=1/3\).

\begin{corollary}[Dihedral Gram determinant]
Let \(G\) be the \(4\times4\) Gram matrix of the outward unit normals:
\begin{equation}
 G_{ii}=1,
 \qquad
 G_{ij}=-\cos\theta_{kl}
 \quad(\{i,j,k,l\}=\{1,2,3,4\}).                           \tag{11.6}
\end{equation}
Then
\begin{equation}
 \detm G=0,
 \qquad
 G(\Delta_1,\Delta_2,\Delta_3,\Delta_4)^T=0.               \tag{11.7}
\end{equation}
\end{corollary}

\begin{proof}
Four vectors in \(\R^3\) are linearly dependent, so their Gram determinant
vanishes.  Equation (11.2), divided by \(2\), is precisely the displayed
weighted null relation.
\end{proof}

\subsection{The tetrahedral sine-product identity}

Let the vertices be denoted \(O,A,B,C\).  All angles in the next theorem are
ordinary planar angles inside the triangular faces.

\begin{theorem}[Tetrahedral sine law]
\label{thm:tetra-sine}
\begin{equation}
 \sin\angle OAB\;\sin\angle OBC\;\sin\angle OCA
 =\sin\angle OAC\;\sin\angle OCB\;\sin\angle OBA.          \tag{11.8}
\end{equation}
Choosing each of the four vertices in turn as \(O\) gives four relations, of
which at most three are independent.
\end{theorem}

\begin{proof}
Apply the ordinary planar sine law in the three faces through \(O\):
\begin{align*}
 \frac{\sin\angle OAB}{\sin\angle OBA}&=\frac{OB}{OA},\\
 \frac{\sin\angle OBC}{\sin\angle OCB}&=\frac{OC}{OB},\\
 \frac{\sin\angle OCA}{\sin\angle OAC}&=\frac{OA}{OC}.
\end{align*}
Multiplying the three equations makes all edge lengths cancel, leaving (11.8).
The planar sine laws can be imported from ProveIt's constructive triangle
package; every denominator is positive in a nondegenerate face.
\end{proof}

Multiplying any three of the four vertex-based equations and canceling the
common positive sine factors yields the fourth.  This explains constructively
why only three constraints are independent.

\subsection{The \(\ipra\) interpretation}

The tetrahedral proof is even more elementary than the spherical one.  Encode
the four vertex coordinates by fixed rapid names.  Cross products give the
unnormalized face-area vectors exactly.  A bounded square-root routine gives
\(2\Delta_i=\norm{N_i}\), with positive lower bounds supplied by
nondegeneracy.  Define \(\cos\theta_{ij}\) by the normalized negative dot
product of the two outward normals.  Then (11.2) is a coordinate polynomial
identity and (11.4) is one ring expansion.  Actual angle names are optional;
the cosine law needs only their cosine values.  Consequently the object-level
proof requires no analytic trigonometry beyond normalization by positive
square roots.
